\chapter{Complexity — Meyer Sigma EXP}
%%% generated by the blueprint pipeline from TCSlib.Complexity.CircuitComplexity.MeyerSigmaEXP
\section{Overview}

Declarations of the Lean module \texttt{TCSlib.Complexity.CircuitComplexity.MeyerSigmaEXP}.

\section{Declarations}

\begin{theorem}[Exact Polynomial Padding Split]
\lean{Complexity.Meyer.solveSplit_self'}
\label{Complexity.Meyer.solveSplit_self'}
\leanok
\uses{Turing.solveSplit}
\difficulty{4}
For all natural numbers $C,c,n$, the least $i\in\{0,\ldots,N\}$ satisfying
$i+C(i+1)^c=N$, with $N=n+C(n+1)^c$, is $n$.
\end{theorem}

\begin{definition}[Encoding a Word by a Valid Split]
\lean{Complexity.Meyer.splitXU}
\label{Complexity.Meyer.splitXU}
\leanok
\uses{Turing.pairEncode, Turing.solveSplit}
For natural numbers $C,c$ and a Boolean word $w$, if the split-solving procedure for
$C,c$, and the length of $w$ returns a split position $i$, define the result to be the
pair encoding of the prefix of $w$ of length $i$ and the suffix after that prefix. If no
split position is returned, define the result to be the empty word.
\end{definition}

\begin{theorem}[Pair Encoding of an Appended String]
\lean{Complexity.Meyer.splitXU_append}
\label{Complexity.Meyer.splitXU_append}
\leanok
\uses{Complexity.Meyer.solveSplit_self', Complexity.Meyer.splitXU, Turing.pairEncode}
\difficulty{1}
Let $C,c$ be natural numbers and let $x,u$ be Boolean strings such that
$|u|=C(|x|+1)^c$. Then splitting the concatenation $x\mathbin{+\!\!+}u$ by the polynomial
padding rule yields the pair encoding of $x$ and $u$: the bits of $x$ are doubled,
followed by the separator $01$ and then $u$.
\end{theorem}

\begin{theorem}[Time-bounded complementation]
\lean{Complexity.Meyer.exists_compl_decider}
\label{Complexity.Meyer.exists_compl_decider}
\leanok
\uses{Turing.FinTM, Turing.FinTM.ComputesFunInTime, Turing.FinTM.DecidesInTime, Turing.FinTM.computesFunInTime_comp, Turing.FinTM.computesFunInTime_ifEq, Turing.MultiTapeTM.indicator}
\difficulty{4}
Let $L$ be a language of finite Boolean words, let $M$ be a finite Turing machine over
the Boolean alphabet, and let $T:\mathbb{N}\to\mathbb{N}$. If $M$ decides $L$ on every
input of length $n$ within $T(n)$ steps, then there exist a finite Boolean-alphabet
Turing machine $M'$ and natural numbers $c,a$ such that $M'$ decides the complement of
$L$ within $c\bigl(T(n)+a(T(n)+1)+1\bigr)$ steps on every input of length $n$.
\end{theorem}

\begin{theorem}[The second existential level lies in exponential time]
\lean{Complexity.SigmaP_two_subset_EXP}
\label{Complexity.SigmaP_two_subset_EXP}
\leanok
\uses{Complexity.EXP, Complexity.ExpPoly, Complexity.ExpPoly.add, Complexity.ExpPoly.comp_poly, Complexity.ExpPoly.mem_EXP, Complexity.ExpPoly.mul, Complexity.ExpPoly.of_le, Complexity.Meyer.exists_compl_decider, Complexity.Meyer.splitXU, Complexity.Meyer.splitXU_append, Complexity.NP_subset_EXP, Complexity.SigmaP, Complexity.exists_proj_decider, Complexity.expPoly_exp, Complexity.expPoly_poly, Complexity.mem_SigmaP_two_iff, Turing.FinTM.ComputesFunInTime, Turing.FinTM.DecidesInTime, Turing.FinTM.computesFunInTime_comp, Turing.FinTM.computesFunInTime_splitSolve, Turing.MultiTapeTM.indicator}
\difficulty{6}
Every language in $\Sigma_2^P$ belongs to
$\mathsf{EXP}=\bigcup_{c\in\mathbb{N}}\mathrm{DTIME}(n\mapsto 2^{n^c})$.
\end{theorem}
